% =====================================================================
% Meridian Financial — DPN Implementation Workshop Booklet
% Compile: pdflatex 2026-08-31-meridian-financial-training-booklet.tex (run twice for refs)
% =====================================================================
\documentclass[10pt,a4paper]{article}

\usepackage[a4paper, margin=1.8cm, top=2.2cm, bottom=2.2cm]{geometry}
\usepackage[utf8]{inputenc}
\usepackage[T1]{fontenc}

\usepackage{xcolor}
\usepackage{booktabs}
\usepackage{array}
\usepackage{tabularx}
\usepackage{multicol}
\usepackage{listings}
\usepackage{fancyhdr}
\usepackage{parskip}
\usepackage{hyperref}

% --- Colors ---
\definecolor{ddpurple}{HTML}{632CA6}
\definecolor{ddlightpurple}{HTML}{C5B0E8}
\definecolor{ddpale}{HTML}{F0EBF9}
\definecolor{ddgreen}{HTML}{00BF9A}
\definecolor{ddgray}{HTML}{6B6B76}
\definecolor{dddark}{HTML}{1E1E2D}
\definecolor{ddorange}{HTML}{F26522}
\definecolor{codebg}{HTML}{F4F4F4}

% --- Hyperref ---
\hypersetup{
  colorlinks=true,
  linkcolor=ddpurple,
  urlcolor=ddpurple,
  pdftitle={Meridian Financial - DPN Implementation Workshop Booklet},
  pdfauthor={Datadog Partner Solutions}
}

% --- Section Styling ---
\renewcommand{\section}[1]{%
  \refstepcounter{section}%
  \addcontentsline{toc}{section}{\protect\numberline{\thesection}#1}%
  \vspace{10pt}%
  \noindent\colorbox{ddpurple}{%
    \parbox{\dimexpr\linewidth-2\fboxsep\relax}{%
      \vspace{3pt}\hspace{4pt}{\large\bfseries\color{white}#1}\vspace{3pt}%
    }%
  }\vspace{6pt}\par%
}

\renewcommand{\subsection}[1]{%
  \vspace{8pt}%
  {\normalsize\bfseries\color{ddpurple}#1}\par\vspace{2pt}%
}

% --- Header / Footer ---
\pagestyle{fancy}
\fancyhf{}
\fancyhead[L]{\small\color{ddgray}\textbf{Meridian Financial -- DPN Implementation Workshop}}
\fancyhead[R]{\small\color{ddgray}August 2026}
\fancyfoot[C]{\small\color{ddgray}\thepage}
\renewcommand{\headrulewidth}{0.4pt}

% --- Box helpers ---
\newcommand{\objectivebox}[1]{%
  \noindent\colorbox{ddpale}{%
    \begin{minipage}{\dimexpr\linewidth-2\fboxsep\relax}%
      \vspace{4pt}%
      {\small\bfseries\color{ddpurple}Learning Objectives}\par\vspace{2pt}%
      #1%
      \vspace{2pt}%
    \end{minipage}%
  }\par\vspace{6pt}%
}

\newcommand{\exercisebox}[1]{%
  \noindent\fcolorbox{ddgreen}{white}{%
    \begin{minipage}{\dimexpr\linewidth-2\fboxrule-2\fboxsep\relax}%
      \vspace{4pt}%
      {\small\bfseries\color{ddgreen}Hands-On Exercise}\par\vspace{2pt}%
      #1%
      \vspace{2pt}%
    \end{minipage}%
  }\par\vspace{6pt}%
}

\newcommand{\tipbox}[1]{%
  \noindent\fcolorbox{ddorange}{white}{%
    \begin{minipage}{\dimexpr\linewidth-2\fboxrule-2\fboxsep\relax}%
      \vspace{3pt}%
      #1%
      \vspace{3pt}%
    \end{minipage}%
  }\par\vspace{6pt}%
}

\newcommand{\resourcebox}[1]{%
  \noindent\fcolorbox{ddgray}{codebg}{%
    \begin{minipage}{\dimexpr\linewidth-2\fboxrule-2\fboxsep\relax}%
      \vspace{4pt}%
      {\small\bfseries\color{ddgray}Datadog Documentation}\par\vspace{2pt}%
      #1%
      \vspace{2pt}%
    \end{minipage}%
  }\par\vspace{6pt}%
}

% --- Code Listings ---
\lstdefinestyle{ddcode}{
  backgroundcolor=\color{codebg},
  basicstyle=\small\ttfamily,
  breaklines=true,
  breakatwhitespace=true,
  frame=single,
  framesep=4pt,
  xleftmargin=4pt,
  xrightmargin=4pt,
  showstringspaces=false,
  tabsize=2,
  aboveskip=4pt,
  belowskip=4pt,
}
\lstset{style=ddcode}

% --- Table helpers ---
\newcolumntype{L}[1]{>{\raggedright\arraybackslash}p{#1}}
\newcolumntype{C}[1]{>{\centering\arraybackslash}p{#1}}

\setlength{\parindent}{0pt}

% =====================================================================
\begin{document}

% =====================================================================
% TITLE PAGE
% =====================================================================
\thispagestyle{empty}
\vspace*{2cm}
\begin{center}

{\fontsize{28}{34}\selectfont\bfseries\color{ddpurple}
DPN Implementation\\[6pt]
Workshop -- Meridian Financial}\\[1.2cm]

{\large\color{dddark} Partner Training -- Attendee Booklet}\\[0.4cm]
{\normalsize\color{ddgray} August 2026}\\[3cm]

\noindent\colorbox{ddpurple}{%
  \begin{minipage}{12cm}%
    \centering\vspace{10pt}
    {\large\bfseries\color{white}9 Modules $\cdot$ make-Pipeline Hands-On Lab}\\[6pt]
    {\normalsize\color{ddlightpurple}Organized around the real \texttt{make} targets in
    \texttt{sample/impl}, not product pillars}\\[10pt]
  \end{minipage}%
}

\end{center}
\vspace{2cm}

\textbf{\color{ddpurple}What changed from the previous booklet:} this curriculum replaces the
earlier product-pillar structure (Foundation / Three Pillars / DEM / Observability / Shift
Left-Right / AI) with modules that map one-to-one onto the actual pipeline stages in the
Meridian Financial sample app's \texttt{Makefile}: \texttt{make tags} $\rightarrow$
\texttt{make dbm} $\rightarrow$ \texttt{make instrument} $\rightarrow$ \texttt{make dem}
$\rightarrow$ \texttt{make security} $\rightarrow$ \texttt{make tf-apply-dd}, plus a
fault-injection diagnosis lab and a dedicated AWS/EKS module. There is no DogStatsD anywhere in
this curriculum -- every custom metric in the sample app is span-based.

\newpage

% =====================================================================
{\Large\bfseries\color{ddpurple}Contents}\par\vspace{8pt}
\makeatletter\@starttoc{toc}\makeatother
\newpage

% =====================================================================
\section{Module 0 -- Project \& Customer Scoping}

Before any \texttt{make} target runs, a Datadog deployment needs a scope. Meridian Financial is
a fictional mid-market financial services company -- payments, account management, nightly
reconciliation, async fraud detection -- used throughout this curriculum.

\objectivebox{
\begin{itemize}\setlength\itemsep{0pt}
\item Read a customer's architecture for monitoring blind spots before proposing anything
\item Map each architecture component to the Datadog capability that closes its gap
\item Translate customer pain points into business impact and priority
\item Produce a phased, justified proposal -- not a feature checklist
\end{itemize}}

\textbf{Duration:} 45--60 minutes \quad \textbf{Prerequisites:} none.

\subsection{The four-step scoping method}
\begin{enumerate}\setlength\itemsep{0pt}
\item \textbf{Read the Architecture} -- critical path, infrastructure beyond the app layer,
  monitoring gaps, ownership.
\item \textbf{Build the Initial Proposal} -- map each component to the Datadog capability that
  closes its blind spot.
\item \textbf{Listen to Customer Requirements} -- map real pain points to business impact.
\item \textbf{Adjust and Prioritize} -- phase the proposal: revenue impact first, compliance/risk
  second, operational efficiency third.
\end{enumerate}

\begin{center}
\begin{tabularx}{\linewidth}{L{4.5cm}L{4.5cm}L{5cm}}
\toprule
\textbf{Architecture component} & \textbf{Blind spot} & \textbf{Datadog capability}\\
\midrule
gateway-api / account-service / transaction-service & No distributed trace across services & APM, Unified Service Tagging\\
PostgreSQL (postgres-ledger) & Query performance invisible without DBA access & Database Monitoring\\
ActiveMQ Artemis queues & Can't tell producer overload from consumer lag & Data Streams Monitoring\\
batch-processor (Spring Batch) & A job can report success while doing less work & Data Jobs Monitoring\\
Finance dashboard (frontend) & No visibility into real user sessions & Real User Monitoring\\
All 6 services & CVEs, injection attacks, container drift & ASM, CWS, CSPM\\
\bottomrule
\end{tabularx}
\end{center}

\exercisebox{Fill in the architecture-to-capability table with one Meridian-specific
justification per row. Pick two customer pain points and assign each a priority phase with one
success criterion. \textit{Expected outcome:} a one-page phased proposal grounded in the real
architecture, mapping onto Modules 1--8.}

\resourcebox{
\begin{itemize}\setlength\itemsep{0pt}
\item Getting Started with Datadog: \url{https://docs.datadoghq.com/getting_started/}
\item Getting Started with Tagging: \url{https://docs.datadoghq.com/getting_started/tagging/}
\end{itemize}}

\newpage
% =====================================================================
\section{Module 1 -- Foundation \& Deploy}

Org setup is general and applies to any customer; deploying Meridian Financial ``cold'' -- with
every Datadog signal off -- is the starting point every later module instruments from.

\objectivebox{
\begin{itemize}\setlength\itemsep{0pt}
\item Configure a new Datadog org: authentication, users, teams, API/App keys
\item Understand when to use single vs. multi-org architecture
\item Deploy the app with zero instrumentation active and confirm it runs cleanly
\item Install the Datadog Agent without enabling any instrumentation layer yet
\end{itemize}}

\textbf{Duration:} 1--1.5 hours \quad \textbf{Prerequisites:} Datadog sandbox org (Admin), a
local cluster or AWS account, \texttt{kubectl}/\texttt{helm}/\texttt{terraform}.

\subsection{Organization setup}
\begin{center}
\begin{tabularx}{\linewidth}{L{4cm}L{5cm}L{5cm}}
\toprule
\textbf{Key type} & \textbf{Purpose} & \textbf{Who holds it}\\
\midrule
API key & Authenticates data sent to Datadog (Agent, Terraform) & The org -- secrets manager\\
App key & Authenticates API calls on a user's behalf & Each individual user\\
\bottomrule
\end{tabularx}
\end{center}

\tipbox{\textbf{Get the App Key \emph{value}}, not the Key ID -- the list view shows the Key ID
by default. Using the Key ID causes \texttt{401 Unauthorized} on every \texttt{terraform apply}.}

\subsection{Deploy Meridian Financial cold}
\begin{lstlisting}
make build            # build all 6 service images + traffic-generator
make deploy-k8s       # deploy the app into the finance namespace
make deploy-k8s-dd    # install Operator, create secret, deploy the Agent
\end{lstlisting}

Confirm: \texttt{kubectl get pods -n finance} (12 Running), log in to the dashboard at
\texttt{http://localhost:30080} as \texttt{carol.admin} / \texttt{Finance@2025!}. At this point:
no traces, no DBM, no RUM -- the deliberate starting state.

\exercisebox{Complete the org setup decisions for your sandbox. Deploy the app and Agent cold.
Confirm APM $>$ Services shows nothing yet (expected). \textit{Expected outcome:} a running
baseline deployment with a documented org configuration.}

\resourcebox{
\begin{itemize}\setlength\itemsep{0pt}
\item Managing Multiple Organizations: \url{https://docs.datadoghq.com/account_management/multi_organization/}
\item SAML SSO: \url{https://docs.datadoghq.com/account_management/saml/}
\item Teams: \url{https://docs.datadoghq.com/account_management/teams/}
\item API \& Application Keys: \url{https://docs.datadoghq.com/account_management/api-app-keys/}
\end{itemize}}

\newpage
% =====================================================================
\section{Module 2 -- Unified Service Tagging (\texttt{make tags})}

\objectivebox{
\begin{itemize}\setlength\itemsep{0pt}
\item Design a tag taxonomy grounded in real business services
\item Apply UST + log injection to all six services and verify correlation end-to-end
\item Recognize each language's log-injection mechanism -- Python, Node, Java, Go all differ
\end{itemize}}

\textbf{Duration:} 45--60 minutes \quad \textbf{Prerequisites:} Module 1 complete.

\begin{center}
\begin{tabularx}{\linewidth}{L{3cm}L{3cm}L{8cm}}
\toprule
\textbf{Language} & \textbf{Service(s)} & \textbf{Log-injection mechanism}\\
\midrule
Python & gateway-api, fraud-detection & \texttt{ddtrace.contrib.logging.patch} + \texttt{patch\_logging()}\\
Node.js & transaction-service & \texttt{logInjection: true} in the \texttt{dd-trace} init block\\
Java & account-service, batch-processor & \texttt{DD\_LOGS\_INJECTION=true} (Logback/Log4j2 MDC hook)\\
Go & notification-service & Manual \texttt{dd.trace\_id}/\texttt{dd.span\_id} injection -- no automatic MDC hook\\
\bottomrule
\end{tabularx}
\end{center}

\tipbox{\textbf{Ordering matters for Go.} \texttt{notification-service}'s log injection reads
the \texttt{alert.send} span's trace/span ID -- a span that only exists once Module 4's
\texttt{make instrument} has uncommented it. Running \texttt{make tags} alone first leaves an
undefined \texttt{span} variable and the build fails.}

\begin{lstlisting}
make tags && make build && make deploy-k8s   # apply UST + log injection, rebuild, redeploy
make untag && make build && make deploy-k8s  # reverse
\end{lstlisting}

\textbf{Validate:} any trace/log carries \texttt{env:staging service:<name> version:latest}. In
Log Explorer, a \textbf{View Trace} button appears on any log once \texttt{dd.trace\_id} is
present.

\exercisebox{Apply UST + log injection. Confirm env vars are set on each pod. Prove the
log$\rightarrow$trace jump works for at least one service, then reverse it.}

\resourcebox{
\begin{itemize}\setlength\itemsep{0pt}
\item Unified Service Tagging: \url{https://docs.datadoghq.com/getting_started/tagging/unified_service_tagging/}
\item Log-Trace Correlation: \url{https://docs.datadoghq.com/tracing/other_telemetry/connect_logs_and_traces/}
\item Tag Policies: \url{https://docs.datadoghq.com/monitors/settings/}
\end{itemize}}

\newpage
% =====================================================================
\section{Module 3 -- Database Monitoring (\texttt{make dbm})}

\objectivebox{
\begin{itemize}\setlength\itemsep{0pt}
\item Understand what DBM collects and why it needs its own database role
\item Apply \texttt{make dbm} and confirm query metrics and explain plans appear
\item Understand the APM $\leftrightarrow$ DBM correlation link
\end{itemize}}

\textbf{Duration:} 30--45 minutes \quad \textbf{Prerequisites:} Module 2 complete.

Two steps: (a) Agent-side config uncomments the \texttt{postgres.d} check + password wiring;
(b) a dedicated read-only \texttt{datadog} role is created in PostgreSQL, granted
\texttt{pg\_stat\_statements} and the \texttt{datadog.explain\_statement} function. The
monitoring role is \textbf{read-only} (\texttt{pg\_monitor} only, never superuser) -- important
for a ledger database specifically.

\begin{lstlisting}
make dbm
kubectl apply -k deploy/kubernetes/datadog/agent
kubectl rollout restart daemonset/datadog-agent -n datadog
make undbm     # reverse: drops the role, reverses the Agent-side patch
\end{lstlisting}

\textbf{Validate:} Databases $\rightarrow$ Query Metrics shows \texttt{postgres-ledger}
queries; open a sample $\rightarrow$ \textbf{Explain Plan}.

\exercisebox{Enable DBM. Generate payment traffic. Open a \texttt{POST /v1/payments} trace,
find the \texttt{ledger.velocity\_check} span, click \textbf{View in DBM}, confirm it lands on
that exact query's explain plan. Reverse it.}

\resourcebox{
\begin{itemize}\setlength\itemsep{0pt}
\item Database Monitoring: \url{https://docs.datadoghq.com/database_monitoring/}
\item DBM -- PostgreSQL self-hosted setup: \url{https://docs.datadoghq.com/database_monitoring/setup_postgres/selfhosted/}
\item DBM + APM correlation: \url{https://docs.datadoghq.com/database_monitoring/connect_dbm_and_apm/}
\end{itemize}}

\newpage
% =====================================================================
\section{Module 4 -- APM, Data Streams \& Data Jobs Monitoring (\texttt{make instrument})}

The densest module -- four steps under one sentinel: APM custom spans, Single Step
Instrumentation gating, Continuous Profiler, and DSM/DJM.

\objectivebox{
\begin{itemize}\setlength\itemsep{0pt}
\item Understand which spans are always-on vs. patch-gated, per service
\item Enable Single Step Instrumentation and verify tracer injection actually happened
\item Correlate a Continuous Profiler flame graph with a slow trace
\item Understand DSM's producer-vs-consumer distinction -- including a real STOMP coverage gap
\item Understand what Data Jobs Monitoring adds on top of a job's own status field
\end{itemize}}

\textbf{Duration:} 1.5--2 hours \quad \textbf{Prerequisites:} Modules 2 and 3.

\textbf{No DogStatsD anywhere.} Every \texttt{finance.*} custom metric is span-based
(\texttt{datadog\_spans\_metric}), generated directly from the spans below.

\subsection{The STOMP / DSM gap}

\tipbox{\texttt{transaction-service} publishes to \texttt{fraud.score.queue} over STOMP via the
\texttt{stompit} library -- and dd-trace Node.js has \textbf{no automatic STOMP instrumentation
at all}. Without extra work it never appears as a producer in the DSM pathway map, no matter how
much traffic it generates. \texttt{make instrument} closes this with a \textbf{manual DSM
checkpoint} (\texttt{setProduceCheckpoint}) in \texttt{producer.js}, plus a matching manual
consume checkpoint in \texttt{fraud-detection/listener.py}. Module 7's diagnosis lab Scenario 3
depends entirely on this checkpoint being in place.}

\begin{center}
\begin{tabularx}{\linewidth}{L{3.5cm}L{3.5cm}L{7cm}}
\toprule
\textbf{Service} & \textbf{Library} & \textbf{Injection mechanism}\\
\midrule
gateway-api, fraud-detection & ddtrace (Python) & \texttt{PYTHONPATH} + auto-instrumentation\\
transaction-service & dd-trace (Node.js) & \texttt{NODE\_OPTIONS=--require dd-trace/init}\\
account-service, batch-processor & dd-java-agent & \texttt{JAVA\_TOOL\_OPTIONS=-javaagent:...}\\
notification-service & dd-trace-go & Not single-step injected -- in-code \texttt{tracer.Start()}\\
\bottomrule
\end{tabularx}
\end{center}

\tipbox{\textbf{fraud-detection needs a rebuild, not just a redeploy.} Its DSM support is a
baked pip dependency, unlike the other services' plain env-var toggle. After
\texttt{make instrument}, run \texttt{make build} (or \texttt{build-ecr}) and reload the image
before DSM data appears for it.}

\begin{lstlisting}
make instrument && make build && make deploy-k8s
make uninstrument && make build   # reverse: DSM/DJM -> profiler -> SSI -> APM spans
\end{lstlisting}

\exercisebox{Enable all four steps. Verify custom spans, SSI injection, profiler flame graphs.
Confirm \texttt{transaction-service} appears as a \textbf{producer} in the DSM pathway map for
\texttt{fraud.score.queue} -- this is the check Module 7 Scenario 3 depends on. Trigger a
reconciliation run and confirm it appears under APM $\rightarrow$ Data Jobs.}

\resourcebox{
\begin{itemize}\setlength\itemsep{0pt}
\item Custom instrumentation: \url{https://docs.datadoghq.com/tracing/trace_collection/custom_instrumentation/}
\item Single-step instrumentation: \url{https://docs.datadoghq.com/tracing/trace_collection/automatic_instrumentation/single-step-apm/}
\item Continuous Profiler: \url{https://docs.datadoghq.com/profiler/}
\item Data Streams Monitoring: \url{https://docs.datadoghq.com/data_streams/}
\item Data Streams -- Node.js manual checkpoints: \url{https://docs.datadoghq.com/data_streams/nodejs/}
\item Data Jobs Monitoring: \url{https://docs.datadoghq.com/data_jobs/}
\end{itemize}}

\newpage
% =====================================================================
\section{Module 5 -- Real User Monitoring (\texttt{make dem})}

\objectivebox{
\begin{itemize}\setlength\itemsep{0pt}
\item Understand how \texttt{make dem} differs mechanically -- a direct API call, not a patch
\item Enable Browser RUM + Session Replay and confirm a session is captured
\item Apply PII masking correctly for a financial application's session replay
\end{itemize}}

\textbf{Duration:} 30--45 minutes \quad \textbf{Prerequisites:} App running (Module 1). Independent
of Modules 2--4 and 7 -- no Terraform dependency.

\texttt{make dem} creates the RUM application via a direct \texttt{POST
/api/v2/rum/applications} call and injects the resulting \texttt{client\_token} into
\texttt{frontend-stub/index.html} with \texttt{sed}. Idempotent via \texttt{.dem-applied}.

\tipbox{\textbf{PII masking matters specifically here} -- this is a financial dashboard.
\texttt{defaultPrivacyLevel: 'mask-user-input'} ensures form values (account numbers, payment
amounts) are never captured in Session Replay. Never disable this on real financial data.}

\begin{lstlisting}
make dem
kubectl create configmap frontend-dashboard --from-file=index.html=frontend-stub/index.html -n finance --dry-run=client -o yaml | kubectl apply -f -
kubectl rollout restart deployment/frontend -n finance
make undem   # reverse
\end{lstlisting}

\textbf{Validate:} RUM $\rightarrow$ Applications $\rightarrow$ \texttt{finance-frontend}
$\rightarrow$ Sessions $\rightarrow$ Session Replay available.

\exercisebox{Enable RUM. Click through the dashboard. Confirm Session Replay shows the
interaction with form fields masked, not literal values. Reverse it.}

\resourcebox{
\begin{itemize}\setlength\itemsep{0pt}
\item Browser RUM: \url{https://docs.datadoghq.com/real_user_monitoring/browser/}
\item RUM Session Replay: \url{https://docs.datadoghq.com/real_user_monitoring/session_replay/}
\item RUM Privacy / PII masking: \url{https://docs.datadoghq.com/real_user_monitoring/session_replay/privacy_options/}
\end{itemize}}

\newpage
% =====================================================================
\section{Module 6 -- Security (\texttt{make security})}

\objectivebox{
\begin{itemize}\setlength\itemsep{0pt}
\item Understand the two-sided nature of \texttt{make security} -- Agent config + app-side flag
\item Distinguish what ASM Threats, ASM SCA, CWS, and CSPM each detect
\item Enable the stage and confirm CWS self-tests pass
\end{itemize}}

\textbf{Duration:} 30--45 minutes \quad \textbf{Prerequisites:} Module 4 -- ASM Threats needs the
tracer already injected via Single Step Instrumentation.

\begin{center}
\begin{tabularx}{\linewidth}{L{3cm}L{3.5cm}L{7cm}}
\toprule
\textbf{Product} & \textbf{Layer} & \textbf{Detects}\\
\midrule
ASM Threats & APM tracer (app-side) & SQLi, XSS, SSRF, credential stuffing, business-logic attacks\\
ASM SCA & Agent-side & Known CVEs in dependencies\\
CWS & Agent eBPF (kernel) & Shell spawned in container, privilege escalation, syscall anomalies\\
CSPM & Agent + cloud APIs & Privileged pods, exposed secrets, CIS/PCI-DSS findings\\
\bottomrule
\end{tabularx}
\end{center}

\begin{lstlisting}
make security
kubectl apply -k deploy/kubernetes/datadog/agent && kubectl rollout restart daemonset/datadog-agent -n datadog
make build && make deploy-k8s   # DD_APPSEC_ENABLED needs a redeploy, not a bare restart
make unsecurity   # reverse
\end{lstlisting}

\textbf{Validate:} \texttt{DD\_APPSEC\_ENABLED=true} on \texttt{gateway-api};
\texttt{security-agent status} shows CWS self-tests succeeded (\texttt{rule\_open},
\texttt{rule\_chmod}, \texttt{rule\_chown}).

\exercisebox{Enable all four products. Confirm both halves independently. Configure one
finance-specific threat rule (e.g. brute force on \texttt{/v1/payments}). Reverse it.}

\resourcebox{
\begin{itemize}\setlength\itemsep{0pt}
\item Application Security (ASM): \url{https://docs.datadoghq.com/security/application_security/}
\item Cloud Workload Security (CWS): \url{https://docs.datadoghq.com/security/cloud_workload_security/}
\item Cloud Security Posture Management (CSPM): \url{https://docs.datadoghq.com/security/cloud_security_management/misconfigurations/}
\end{itemize}}

\newpage
% =====================================================================
\section{Module 7 -- Dashboards, Monitors \& the Diagnosis Lab (\texttt{make tf-apply-dd})}

\objectivebox{
\begin{itemize}\setlength\itemsep{0pt}
\item Apply the Terraform-managed dashboard, monitors, SLOs, and synthetic tests
\item Diagnose three realistic incidents live, using the exact capability that catches each one
\item Understand why a monitor's absolute thresholds can fail to fire at demo scale even when
  the underlying signal is genuinely correct
\end{itemize}}

\textbf{Duration:} 1.5--2 hours \quad \textbf{Prerequisites:} Modules 2--4.

\begin{lstlisting}
make tf-apply-dd     # DD_API_KEY / DD_APP_KEY read from .env automatically; dashboard,
                     # 7+ monitors, 3 SLOs, 7 synthetic tests, log pipeline
make tf-destroy-dd    # reverse -- WARNING: deletes the log index and indexed logs too
\end{lstlisting}

\subsection{Scenario 1 -- Payments Are Slow and No One Knows Why}

\texttt{make scenario-1} drops \texttt{idx\_transactions\_account\_id}, turning the
\texttt{ledger.velocity\_check} query into a full table scan on every payment. The
\texttt{transactions} table is pre-seeded with $\sim$300k backdated rows specifically so this
scan is slow enough to see -- on live traffic volume alone it finishes in under a millisecond.
\textbf{Capability:} APM traces + log correlation + the APM $\leftrightarrow$ DBM link.

\subsection{Scenario 2 -- The Nightly Reconciliation Is Failing Silently}

\texttt{make scenario-2} appends \texttt{AND currency <> 'JPY'} to the reconciliation reader's
\texttt{WHERE} clause -- the job still completes with no error; only the record count drops
$\sim$20\%. \textbf{Capability:} Data Jobs Monitoring.

\tipbox{The \texttt{reconciliation\_low\_record\_count} monitor's absolute thresholds were
sized for production volume -- at workshop traffic scale, the $\sim$20\% drop may not cross
them within the session. \texttt{job.records\_processed} in Data Jobs Monitoring is the live,
always-reliable signal; treat the monitor as ``here's the kind of alert you'd build,'' not a
guaranteed live-fire.}

\subsection{Scenario 3 -- The Fraud Queue Is Backing Up}

\texttt{make scenario-3} triples \texttt{transaction-service}'s publish rate to
\texttt{fraud.score.queue}. The producer floods the queue; \texttt{fraud-detection} keeps up at
its normal rate -- scaling the consumer would do nothing. \textbf{Depends on Module 4's DSM
manual checkpoint} -- without it \texttt{transaction-service} never appears as a producer at all.
\textbf{Capability:} Data Streams Monitoring pathway map + Deployment Tracking.

\tipbox{The \texttt{fraud\_queue\_depth} monitor may \textbf{not} fire at workshop scale --
\texttt{fraud-detection} keeps up even at 3x load, so depth stays near 0. The
\texttt{fraud\_queue\_producer\_surge} monitor (a change alert on the production
\emph{rate}, not depth) is the one that actually catches this scenario live.}

\exercisebox{Run all three scenarios in turn. Diagnose each using only the intended capability,
before checking this booklet's answer. Note explicitly which monitors fired and which didn't,
and why. Reverse all three, then \texttt{tf-destroy-dd} if ending the session.}

\resourcebox{
\begin{itemize}\setlength\itemsep{0pt}
\item Synthetic Monitoring: \url{https://docs.datadoghq.com/synthetics/}
\item Synthetic $\rightarrow$ APM correlation: \url{https://docs.datadoghq.com/synthetics/apm/}
\item Generate metrics from spans: \url{https://docs.datadoghq.com/tracing/trace_pipeline/generate_metrics/}
\end{itemize}}

\newpage
% =====================================================================
\section{Module 8 -- AWS/EKS Deployment}

\objectivebox{
\begin{itemize}\setlength\itemsep{0pt}
\item Provision EKS/VPC/ECR/IAM via Terraform and confirm the cluster is reachable
\item Understand the two AWS-specific integration points this app needs
\item Tear down AWS resources in the correct order
\end{itemize}}

\textbf{Duration:} 1--1.5 hours \quad \textbf{Prerequisites:} AWS CLI $\geq$ 2.x, an SSO
profile, Terraform $\geq$ 1.5.

\tipbox{\textbf{RUM validation status:} Module 5's Browser RUM was validated locally in this
curriculum's development. \textbf{It has not been validated end-to-end on AWS EKS} -- treat it
as untested territory on this path, and verify it yourself before relying on it with a customer.}

\begin{lstlisting}
aws sso login --profile <profile>
make tf-plan-aws && make tf-apply-aws     # provisions EKS/VPC/ECR/IAM/NLB, ~15-20 min
make tf-configure-kubectl && kubectl get nodes
make build-ecr && make deploy-k8s-eks
\end{lstlisting}

Two AWS-specific integration points Terraform doesn't own: Keycloak's public URL (patch the
\texttt{app-config} ConfigMap and the frontend HTML to point at the NLB) and Keycloak's allowed
CORS origins (the realm import hardcodes \texttt{localhost} -- patch via \texttt{kcadm.sh} after
the NLB hostname is known). Both are detailed in Module 8's full write-up
(\texttt{docs/Module-8-AWS-EKS-Deployment.md}).

\tipbox{Both patches are \textbf{in-memory only} -- Keycloak runs \texttt{start-dev} with an
in-process H2 database, so any pod restart wipes them back to \texttt{localhost}-only. If login
fails with \texttt{Failed to fetch} again after working before, re-apply the CORS patch.}

\begin{lstlisting}
make deploy-k8s-dd
make tf-apply-dd
# Teardown, in this order:
make tf-destroy-dd
make tf-destroy-aws
\end{lstlisting}

\exercisebox{Provision, deploy, patch Keycloak's URL and CORS origins, log in via the NLB URL,
apply Datadog resources, then tear down completely. \textit{Expected outcome:} a working EKS
deployment confirmed end-to-end, followed by a clean teardown with no orphaned AWS resources.}

\resourcebox{
\begin{itemize}\setlength\itemsep{0pt}
\item AWS Integration: \url{https://docs.datadoghq.com/integrations/amazon_web_services/}
\item Datadog on Kubernetes (Operator/Helm, Cluster Agent): \url{https://docs.datadoghq.com/containers/kubernetes/}
\item Datadog Operator: \url{https://github.com/DataDog/datadog-operator}
\item Terraform provider: \url{https://registry.terraform.io/providers/DataDog/datadog/latest/docs}
\end{itemize}}

\vspace{1cm}
\begin{center}
{\small\color{ddgray} Source of truth for all modules:
\texttt{sample/impl/INSTRUMENTATION.md}, \texttt{USECASES.md}, \texttt{README.md},
\texttt{TROUBLESHOOTING.md} -- this booklet summarizes, the repo governs.}
\end{center}

\end{document}
